\section{An interior contour removes the ideal return's Gevrey restriction}
\label{sec:interior-contour-source-return}

The reciprocal on a fixed boundary line can have the
logarithmic-square loss of MA-RS-001. The actual theta ideal
has additional structure: its Mellin numerator cancels the
entire nontrivial divisor. An input contour can therefore
move inside the open strip of its inherited two-sided norm.
We construct such a contour with polynomial reciprocal bounds
and use it in the original local-source formula. This
removes both the Gevrey and generic-width assumptions for
that restricted source return. It does not remove them by
an estimate on the unrestricted one-sided inverse.

Retain the nontrivial primitive finite-order character
\(\eta\), structural \(S\), actual compact-sector source
\(\mathcal T_\eta\), Fourier pair, and theta embedding of
MA-MO-012 and MA-MO-017. Put
\[
 \begin{gathered}
 0<a<1/2,\quad s\ge0,\quad
 H=H_{\Sigma,a,s},\quad
 J=\overline{I_\eta^c}^{\,H},\\
 A_\eta(z)=E_{\eta,S}(1/2+z)L_\eta(1/2+z),\\
 \phi=1\text{ sufficiently far left},\quad
 \phi=0\text{ sufficiently far right},\quad
 \phi'\in C_c^\infty,\quad
 \psi=1-\phi-\phi^\vee .
 \end{gathered}
 \tag{MA-MO.800}
\]
The cutoff is real and smooth. Its derivative and the
test \(g\in C_c^\infty(\mathbb R)\) are arbitrary compact
smooth functions; no Gevrey bounds are assumed.
The local multiplier \(m_\eta\), source endomorphism
\(\mathcal L_\eta\), and exact normalizations are those
of (MA-MO.281)--(MA-MO.283).

\paragraph{MA-MO-028 (all-width, all-smooth actual ideal return).}
For every width and polynomial weight in (MA-MO.800)
there is a unique continuous nuclear map
\[
 \begin{aligned}
 \mathcal A_{\eta,g}:J&\longrightarrow\mathcal T_\eta,\\
 \mathcal A_{\eta,g}\vartheta_\eta F
       &=C_g\mathcal L_\eta F
                         \qquad(F\in\mathcal T_\eta).
 \end{aligned}
 \tag{MA-MO.801}
\]
Its theta composition is the actual smoothed common-sum
action on the entire smooth theta ideal and is trace class:
\[
 \mathcal B_{\eta,g}=\vartheta_\eta\mathcal A_{\eta,g}
       \in\mathcal S_1(J),\qquad
 \mathcal B_{\eta,g}|_{I_\eta^c}
                    =B^{\rm raw}_{\eta,g}.
 \tag{MA-MO.802}
\]
The return in fact takes values in the original smooth
source, and its theta composition takes values in
\(I_\eta^c\) with its \(\mathcal E\) topology.
On each fixed compact test-support space,
\(g\mapsto\mathcal B_{\eta,g}\) is continuous into trace
class. The construction and its actual trace comparison
give
\[
 \operatorname{Tr}_{J}\mathcal B_{\eta,g}
 =\frac1{2\pi}\int_{\mathbb R}
      \mathcal Mg(it)\frac{m_\eta'(it)}{m_\eta(it)}\,dt
                  +g(0)\int_{\mathbb R}\psi(x)\,dx .
 \tag{MA-MO.803}
\]
There is no zero-location or simplicity hypothesis,
and a width whose boundary contains zeros is allowed.
This is an assertion on \(J\), not the full Hilbert
completion of the original common-sum source.

\begin{proof}
\emph{Polynomial good paths strictly inside the input strip.}
We first construct a path depending only on the fixed
function and width, not on \(g\) or the source vector.
Use the horizontal heights of MT-M18,
\[
 n+\tfrac14<h_n<n+\tfrac34,\qquad
             \tfrac12<h_{n+1}-h_n<\tfrac32 .
 \tag{MA-MO.804}
\]
On each horizontal segment
\(\{c+ih_n:|c|\le a/4\}\), (M18.8), after canceling
the common Stirling factor, gives
\[
             |A_\eta(c+ih_n)|^{-1}
                     \le C(2+|n|)^D .
 \tag{MA-MO.805}
\]
The finite deleted Euler factors are bounded away from
zero on this fixed positive finite-\(L\) strip.
The height choice avoids every nontrivial zero ordinate.
There are no trivial finite-\(L\) zeros in this strip.
Thus the finitely many bounded-height segments are also
covered after increasing \(C\).

Here is the corresponding vertical choice. For each
sufficiently large \(|n|\), center a disk at
\(2+i(h_n+h_{n+1})/2\), in the uncentered finite-\(L\)
variable, and choose a zero-free boundary radius
\(12<R_n<13\). At these heights its only zeros are
nontrivial zeros, and its total multiplicity \(M_n\)
is \(O(\log(2+|n|))\). For \(c\in[-a/4,a/4]\), put
\[
 V_n(c)=
 \sum_{\rho\text{ in this disk}}
       m_\rho\log^+\frac1{|1/2+c-\Re\rho|}.
 \tag{MA-MO.806}
\]
Each summand has integral at most
\(\int_{-1}^{1}-\log|u|\,du=2\) over this interval.
Since the interval has length \(a/2\), averaging
chooses \(c_n\in[-a/4,a/4]\), avoiding the finitely
many zero real parts in the sum, with
\[
 \begin{aligned}
 V_n(c_n)&\le (4/a+1)M_n,\\
 \prod_{\rho\text{ in the disk}}
       \min\{1,|1/2+c_n-\Re\rho|\}^{m_\rho}
           &\ge e^{-(4/a+1)M_n}
            \ge C^{-1}(2+|n|)^{-D}.
 \end{aligned}
 \tag{MA-MO.807}
\]
When \(M_n=0\), take the midpoint.

We spell out why this product controls the whole
vertical segment, rather than one selected point.
Write \(L_\eta=H_nB_n\) on the disk, where \(B_n\)
is its finite disk Blaschke product, with multiplicities.
The polynomial strip upper bound gives
\(|H_n|\le C(2+|n|)^C\) on the disk. At its center,
the Euler-product lower bound and \(|B_n|\le1\)
give \(|H_n|\ge c>0\).
Harnack's inequality applied to
\(\log(C(2+|n|)^C)-\log|H_n|\) therefore gives
\[
                     |H_n(w)|\ge C^{-1}(2+|n|)^{-D}
 \tag{MA-MO.808}
\]
on the inner half disk.
Every point \(w=1/2+c_n+it\),
\(h_n\le t\le h_{n+1}\), lies in its inner quarter.
Each normalized Blaschke factor there has modulus
at least
\[
       \frac{\min\{1,|1/2+c_n-\Re\rho|\}}{2R_n}.
\]
Multiplying these bounds, using (MA-MO.807) and
\((2R_n)^{M_n}\le C(2+|n|)^D\), proves
\[
 |A_\eta(c_n+it)|^{-1}\le C(2+|n|)^D
                      \qquad(h_n\le t\le h_{n+1}).
 \tag{MA-MO.809}
\]
This is the same finite-disk argument as MT-M18,
but with averaging in the real coordinate instead
of the imaginary coordinate. For the finitely many
remaining blocks, choose \(c_n\) avoiding all zero
real parts meeting the compact vertical rectangle.
Compactness then supplies (MA-MO.809) with an enlarged
constant. No zero-separation hypothesis was used.

Let \(\Gamma\) be the upward-oriented polygonal path
consisting of the vertical segment at \(c_n\) from
height \(h_n\) to \(h_{n+1}\), followed by the
horizontal segment to \(c_{n+1}+ih_{n+1}\).
Its horizontal segments may point in either direction.
Equations (MA-MO.805) and (MA-MO.809) give
\[
 \begin{gathered}
 |\Re w|\le a/4\quad(w\in\Gamma),\\
 \operatorname{length}\bigl(\Gamma
          \cap\{h_n\le\Im w\le h_{n+1}\}\bigr)
                    \le 3/2+a,\\
                  |A_\eta(w)|^{-1}
                       \le C(2+|\Im w|)^D\quad(w\in\Gamma).
 \end{gathered}
 \tag{MA-MO.810}
\]
The length bound permits either convention for assigning
the horizontal endpoints to their adjacent blocks.
The path may cross the central vertical line; this
does not affect any contour identity below.

\emph{The Hilbert norm controls interior evaluations.}
For every \(y\in H\), its Mellin transform is holomorphic
on \(|\Re w|<a\). For \(|\Re w|\le a/4\),
weighted Cauchy--Schwarz gives, for every integer \(j\ge0\),
\[
 \begin{aligned}
 |(\mathcal My)^{(j)}(w)|
 &\le \|y\|_H
 \left(\int_{\mathbb R}
       \frac{|x|^{2j}e^{2\Re(w)x}}
       {(e^{2ax}+e^{-2ax})(1+x^2)^{s/2}}\,dx\right)^{1/2}\\
 &\le C_{a,s,j}\|y\|_H .
 \end{aligned}
 \tag{MA-MO.811}
\]
The last integrand has a fixed exponential decrease
at both ends, uniformly in that interior interval.
Holomorphy and its differentiated formula follow by
the same domination on every smaller open strip.
In particular (MA-MO.810)--(MA-MO.811) imply
\[
       \left|\frac{\mathcal My(w)}{A_\eta(w)}\right|
           \le C(2+|\Im w|)^D\|y\|_H
                                      \qquad(w\in\Gamma).
 \tag{MA-MO.812}
\]
This estimate is for evaluations on the chosen path.
It makes no assertion that the reciprocal is a
Schwartz multiplier on either boundary line.

\emph{The original local kernel as a holomorphic contour kernel.}
Set
\[
 \begin{gathered}
 G(z)=\mathcal Mg(z),\qquad
 P(\xi)=\int e^{\xi x}\phi'(x)\,dx,\qquad
 \Psi(\xi)=\int e^{\xi x}\psi(x)\,dx,\\
 \mathcal K_\eta(z,w)=G(z)\left[
   -\frac{P(z-w)}{z-w}
          \left(\frac{m_\eta(z)}{m_\eta(w)}-1\right)
       +\frac{m_\eta(z)}{m_\eta(w)}\Psi(z-w)\right].
 \end{gathered}
 \tag{MA-MO.813}
\]
For fixed \(\Re z>-1/2\), this is holomorphic in
\(|\Re w|<a\). The apparent singularity \(w=z\)
is removable, since \(m_\eta\) is holomorphic and
nonzero on the possible input strip.
The local multiplier and its derivatives have
polynomial bounds on fixed strips by (MA-MO.281).
Its inverse has polynomial bounds uniformly on
\(|\Re w|\le a/4\), which is strictly inside
\((-1/2,1/2)\).

For every fixed output line \(z=c+it\), \(c>-1/2\),
every derivative order \(j\), and every \(N\), these
bounds and integration by parts in the three compact
smooth transforms give a bound of the form
\[
 \left|\partial_t^j\mathcal K_\eta(c+it,w)\right|
 \le C_{c,j,N}
   \sum_{\nu\le j}|G^{(\nu)}(c+it)|
       (1+|t|+|\Im w|)^{D_{c,j}}
       (1+|t-\Im w|)^{-N},
       \qquad w\in\Gamma.
 \tag{MA-MO.814}
\]
For \(|z-w|\ge1\), this follows directly from
(MA-MO.813). For \(|z-w|\le1\), write
\((m_\eta(z)-m_\eta(w))/(z-w)\) as the integral of
\(m_\eta'\) along the segment. The segment stays in a
fixed closed substrip of \(\Re z>-1/2\). The same
argument applies after every output derivative.
Thus no singular denominator is hidden in (MA-MO.814).
All constants are uniform when \(c\) lies in a fixed
compact interval in this half-plane.

For an actual smooth ideal vector \(y=\vartheta_\eta F\),
the unconditional division statement (MA-MO.173)--(MA-MO.178)
gives
\[
 Q_y(w)=\frac{\mathcal My(w)}{A_\eta(w)}
             =\mathcal MF(w)
       \quad\text{holomorphic for }\Re w>-1/2,
 \tag{MA-MO.815}
\]
with rapid decrease, including derivatives, on every
fixed closed substrip. This is its original source
inverse, not an arbitrary section of the Hilbert vector.
The local formula (MA-MO.288), whose proof uses only
the polynomial local multiplier estimates and compact
smoothness, yields
\[
 \mathcal M(C_g\mathcal L_\eta F)(z)
       =\frac1{2\pi i}\int_{i\mathbb R}
                     \mathcal K_\eta(z,w)Q_y(w)\,dw .
 \tag{MA-MO.816}
\]
The factor \(1/(2\pi i)\) corresponds to the
unnormalized Mellin transforms; on \(w=iu\), this
is exactly the \(1/(2\pi)\) Fourier kernel in
(MA-MO.288).

\emph{Deformation on the actual ideal, before taking its closure.}
Truncate the central contour at two selected heights
and the polygonal contour at the same heights, joining
their endpoints horizontally inside \(|\Re w|\le a/4\).
For the actual smooth ideal input, the integrand in
(MA-MO.816) is holomorphic throughout this strip by
(MA-MO.815), including all its canceled global zeros.
Cauchy's theorem applies to the finite polygon. If
the polygon crosses the central line, apply it to
the successive rectangles; the intermediate sides
cancel with their orientations.
The two final horizontal connectors tend to zero:
their lengths are uniformly bounded, \(Q_y\) is
rapid there, and the local kernel and its derivatives
are polynomially bounded for fixed \(z\).
Equations (MA-MO.810) and (MA-MO.814) give absolute
convergence of the infinite path integral. Consequently
\[
 \mathcal M(C_g\mathcal L_\eta F)(z)
       =\frac1{2\pi i}\int_{\Gamma}
          \mathcal K_\eta(z,w)
               \frac{\mathcal My(w)}{A_\eta(w)}\,dw .
 \tag{MA-MO.817}
\]
There is no residue term: this deformation is performed
only after the actual ideal's full zero jets have
canceled every pole in the quotient.

\emph{Uniform source estimates from the inherited Hilbert norm.}
Combine (MA-MO.810), (MA-MO.812) and
(MA-MO.814). The length per height block is bounded.
For \(N\) sufficiently large compared with the fixed
polynomial exponents,
\[
 \sum_{n\in\mathbb Z}
  (1+|t|+|n|)^{D_1}(1+|n|)^{D_2}
        (1+|t-n|)^{-N}
                \le C(1+|t|)^{D_3}.
 \tag{MA-MO.818}
\]
For example, bound \(1+|n|\) by
\((1+|t|)(1+|t-n|)\), do the same for
\(1+|t|+|n|\), and sum a remaining power
less than \(-2\).
The rapid decrease of every \(G^{(\nu)}(c+it)\)
now gives, for all integers \(M,j\ge0\),
\[
 \sup_{t\in\mathbb R}(1+|t|)^M
  \left|\partial_t^j
        \mathcal M(C_g\mathcal L_\eta F)(c+it)\right|
          \le C_{c,M,j,g,\phi}\|y\|_H .
 \tag{MA-MO.819}
\]
On a fixed compact support for \(g\), each constant
is bounded by finitely many ordinary derivative
seminorms of \(g\). The same statement holds locally
uniformly in \(c>-1/2\).
Fourier inversion and the usual Schwartz estimates
therefore control every seminorm of
\(e^{cx}C_g\mathcal L_\eta F\) by \(\|y\|_H\).
No Gevrey estimate occurs in this calculation.

The actual dual Fourier identity is (MA-MO.295):
\[
 \mathfrak S_\eta C_g\mathcal L_\eta F
      =C_{g^\vee}\mathcal L_{\bar\eta}\mathfrak S_\eta F,
 \qquad
       \vartheta_{\bar\eta}\mathfrak S_\eta F=Ry .
 \tag{MA-MO.820}
\]
Construct an interior path for the dual character
by the same procedure, and apply (MA-MO.819) to
\(Ry\). Reflection is isometric on the exact sum
norm. Hence every weighted Schwartz seminorm of
the actual dual output is also bounded by \(\|y\|_H\).
All conductor, different and compact-character phases
remain those of (MA-MO.820); no unshifted archimedean
Fourier map is substituted.

\emph{Extension to the whole completed ideal and the actual target.}
For a Cauchy sequence from \(I_\eta^c\) in \(H\),
the two output families just estimated converge in
all their weighted Schwartz seminorms. On the central
local \(L^2\) line the actual source Fourier map is
unitary, so their limits remain the correct Fourier
pair. Meyer's fixed-sector source topology is precisely
this paired weighted-Schwartz topology. Thus the
limit belongs to \(\mathcal T_\eta\), continuously
and independently of the approximating sequence.
This defines \(\mathcal A_{\eta,g}\) on all of \(J\)
and proves (MA-MO.801), rather than only a formula
for individually smooth \(y\).

The scalar path integral can be formed on arbitrary
vectors of \(H\), since (MA-MO.811) holds there.
We do not assert the requisite dual-source compatibility
or the original-source identification on those vectors.
The preceding extension uses the actual ideal and
its density. This domain distinction is essential.

\emph{Nuclearity with one bounded family of genuine source vectors.}
We include the functional-analytic step so that the
all-weight return is stronger than separate Hilbert
kernel estimates. Choose increasing Hilbertian norms
\(q_n\) defining the nuclear Fr\'echet topology of
\(\mathcal T_\eta\), each dominating its injective
central local \(L^2\) norm, and sufficiently far apart
that their successive completion maps are trace class.
The source property is the one used and explicitly
located in MA-MO-017. Estimate (MA-MO.819) and its
dual give a bounded map \(J\to(\mathcal T_\eta,q_n)\)
for every \(n\).
Factoring through a stronger norm shows that
\(\mathcal A_{\eta,g}\) is trace class into each
of these Hilbert completions.

Let \(P_N\) be increasing finite-rank orthogonal
projections on \(J\), tending strongly to one.
For each \(n\), the \(q_n\) trace norm of
\(\mathcal A_{\eta,g}(1-P_N)\) tends to zero.
Select an increasing sequence \(N_k\) so that this
tail is at most \(2^{-k}\) for all \(n\le k\).
The singular-value decomposition of
\(\mathcal A_{\eta,g}(P_{N_{k+1}}-P_{N_k})\),
measured in \(q_k\), has coefficient sum at most
\(2^{-k}\), norm-one domain functionals, and
normalized output vectors in the actual source.
They lie in that source because each finite-dimensional
image is spanned by images of vectors of \(J\).
For each fixed \(q_n\), these normalized vectors
are bounded for \(k\ge n\), and the finitely many
earlier blocks are bounded as well. Including the
first finite block gives
\[
 \mathcal A_{\eta,g}y
     =\sum_\nu \lambda_\nu\,\ell_\nu(y)F_\nu,
 \quad
 \sum_\nu|\lambda_\nu|<\infty,\quad
 \sup_\nu\|\ell_\nu\|_{J^*}<\infty,\quad
 \{F_\nu\}\text{ bounded in }\mathcal T_\eta .
 \tag{MA-MO.821}
\]
This is a single genuine nuclear representation.
It is not a collection of unrelated one-sided traces.

The original theta embedding is continuous
\(\mathcal T_\eta\to I_\eta^c\subset\mathcal E\)
by the structural source comparison and nontrivial-sector
Poisson summation in MA-MO-012. Thus its composition
with (MA-MO.821) is trace class on \(J\) and actually
smooths into \(I_\eta^c\). On the dense smooth ideal,
the unconditional source identity (MA-MO.183) gives
\(\vartheta_\eta C_g\mathcal L_\eta F
=B^{\rm raw}_{\eta,g}\vartheta_\eta F\).
Continuity proves (MA-MO.802) everywhere.
To verify test continuity, choose a source Hilbert
norm that controls the theta map into \(H\), then
a stronger norm with trace-class inclusion. Its
operator bound in (MA-MO.819) uses only finitely
many derivative seminorms of \(g\). This bounds
\(\|\mathcal B_{\eta,g}\|_1\) by those same finite
seminorms on each fixed compact test support.

\emph{The actual trace comparison.}
The reverse composition on the original source is
\[
 \mathcal A_{\eta,g}\vartheta_\eta
                       =C_g\mathcal L_\eta .
 \tag{MA-MO.822}
\]
The source has the approximation property from its
split weighted-Schwartz model. Applying the same
absolutely summable pairings in (MA-MO.821) to
the two compositions gives
\[
 \operatorname{Tr}_{J}
          (\vartheta_\eta\mathcal A_{\eta,g})
     =\operatorname{Tr}_{\mathcal T_\eta}
          (\mathcal A_{\eta,g}\vartheta_\eta)
     =\operatorname{Tr}_{\mathcal T_\eta}C_g\mathcal L_\eta .
 \tag{MA-MO.823}
\]
The latter is the local source trace already computed
without a global inverse. Indeed (MA-MO.813) on the
central local line, without \(A_\eta^{-1}\), is a
Schwartz kernel divided by \(2\pi\). Its diagonal is
\[
       \frac{G(it)}{2\pi}
          \left(\frac{m_\eta'(it)}{m_\eta(it)}
                         +\int\psi\right),
 \tag{MA-MO.824}
\]
because \(P(0)=\int\phi'=-1\).
The local-to-source nuclear factorization with global
inverse replaced by one is exactly the unconditional
local argument in MA-MO-017; its paired weighted
estimates use polynomial local multipliers and compact
smooth data only. It equates the actual local \(L^2\)
trace to the source trace. Integrating (MA-MO.824)
and using \((2\pi)^{-1}\int G(it)\,dt=g(0)\)
proves (MA-MO.803).
This concludes the all-width, all-smooth theorem.
\end{proof}

\paragraph{The rational positive source and its two actual moments.}
There are corresponding unconditional, all-width conclusions
for the rational source with \(S=\{\infty\}\).
Use its exact spaces \(\mathscr U_0,\mathscr U_{\rm pf}\),
\(J_{a,s},K_{a,s,\phi}\), and \(\mathcal W_+\) from
MA-MO-010/011/016. The original additive cutoff is
\(v(t)=\phi(\log|t|)\), extended by one at zero.
For \(y=\Theta f\), \(f\in\mathscr U_0\), put
\[
 F(x)=2e^{x/2}f(e^x),\qquad
 \widetilde F(x)=2e^{x/2}\mathfrak Ff(e^x),\qquad
                A(z)=\zeta(1/2+z).
 \tag{MA-MO.825}
\]
The unconditional actual source identities
(MA-MO.127)--(MA-MO.128) identify their Mellin transforms
with \(\mathcal My/A\) and \(\mathcal M(Ry)/A\)
throughout the interior contour strip. Their zero
singularities are canceled, and their weighted transforms
are rapid there. The pole of \(A\) at \(z=1/2\) lies
outside that strip and creates no contour term.
The same path construction therefore applies to both
actual radial sources.

We explicitly check the source moments as continuous
functionals on the completed ideal, rather than importing
an unstated inverse-domain bound. Define
\[
 \begin{aligned}
 C_\phi(w)
   &=\int_{\mathbb R}\phi(x)e^{(1/2-w)x}\,dx
       =-\frac{P(1/2-w)}{1/2-w},\qquad \Re w<1/2,\\
 \lambda_0(y)&=\int_{\mathbb R}v(t)f(t)\,dt,\\
 \lambda_1(y)&=\int_{\mathbb R}v(t)\mathfrak Ff(t)\,dt .
 \end{aligned}
 \tag{MA-MO.826}
\]
The integration by parts for \(C_\phi\) has no endpoint
term: at the left the exponential decreases, and at
the right \(\phi\) is zero. This transform is rapidly
decreasing on every closed vertical substrip of
\(\Re w<1/2\).
Changing the additive variable \(t=e^x\), using
evenness and the factor two in (MA-MO.825), gives
\[
 \begin{aligned}
 \lambda_0(y)
   &=\int_{\mathbb R}\phi(x)e^{x/2}F(x)\,dx\\
   &=\frac1{2\pi i}\int_{\Gamma}
           C_\phi(w)\frac{\mathcal My(w)}{A(w)}\,dw,\\
 \lambda_1(y)
   &=\frac1{2\pi i}\int_{\Gamma}
           C_\phi(w)\frac{\mathcal M(Ry)(w)}{A(w)}\,dw .
 \end{aligned}
 \tag{MA-MO.827}
\]
First derive these identities on the central line by
absolute Fourier inversion and then deform exactly
as in (MA-MO.817). On the path the rapid factor
\(C_\phi\) defeats the polynomial bound (MA-MO.812).
Thus
\[
       |\lambda_0(y)|+|\lambda_1(y)|
                         \le C_{a,s,\phi}\|y\|_{H_{\Sigma,a,s}} .
 \tag{MA-MO.828}
\]
They extend uniquely to all of \(J_{a,s}\).
On overlap with MA-MO-011's regimes, these are the
same actual moment extensions, since their values
agree on the dense smooth theta ideal. Its closure
and codimension proof uses only moment continuity
and two fixed smooth correcting vectors. The
correction (MA-MO.147) therefore repeats at every
width here, and gives
\(K_{a,s,\phi}=J_{a,s}\cap\ker\lambda_0\cap\ker\lambda_1\).
The boundary-line formula involving \(\beta_a\)
in (MA-MO.148) is not asserted when \(j_a\)
is singular. The newly continuous functionals
on \(J_{a,s}\) nevertheless have their canonical
Riesz vectors and the same rank-two Gram construction.

The local kernel (MA-MO.813), with the rational
multiplier (MA-MO.282), and the two radial sources
in (MA-MO.825) now give a continuous nuclear map
\(\mathcal A_{0,g}:J_{a,s}\to\mathscr S_{\rm even}(\mathbb R)\)
agreeing with the actual \(\mathcal D_gLf\) on
\(y=\Theta f\).
Here \(\mathcal D_g\) is the original additive dilation
operator corresponding to centered \(C_g\).
To verify the target, the bounds for all \(c>-1/2\)
give integrability at additive zero and rapid decrease
at infinity; the paired Fourier bounds give integrability
of every Fourier moment. Fourier inversion then gives
all ordinary Schwartz derivatives, including at zero,
exactly as in the rational target verification of
MA-MO-017. The common nuclear-series proof applies to
the actual even Schwartz space.

Add the two actual polar source sections \(f_0,f_1\)
of MA-MO-016:
\[
 \begin{aligned}
 \mathcal A_{+,g}: \mathcal W_+=J_{a,s}\oplus\mathbb C^2
                      &\longrightarrow\mathscr S_{\rm even}(\mathbb R),\\
 \mathcal A_{+,g}(y,c_0,c_1)
       &=\mathcal A_{0,g}y+
                 \sum_{j=0}^1 c_j\mathcal D_gLf_j .
 \end{aligned}
 \tag{MA-MO.829}
\]
The added columns are genuine fixed Schwartz vectors.
Hence this is a nuclear return for every compact smooth
test and every width, with its prescribed source-section
values. The endomorphism
\(\jmath_+\mathcal A_{+,g}\) of \(\mathcal W_+\)
has the local trace (MA-MO.803) with the rational
multiplier, by the same nuclear pairing as above.
It remains distinct from the full original positive
column, whose outgoing polar terms are those in
(MA-MO.299).

Finally \(Lf(0)=-\lambda_1(\Theta f)\) and
\(\int Lf=-\lambda_0(\Theta f)\) on \(\mathscr U_0\),
by the actual source algebra MA-MO-008.
After compact dilation these values are multiplied
by \(G(-1/2)\) and \(G(1/2)\), respectively.
They therefore vanish on \(\mathscr U_{\rm pf}\),
and the continuous extensions (MA-MO.828) show the
same for every input from \(K_{a,s,\phi}\).
The zero-moment source \(\mathscr U_0\) is a closed
finite-codimensional Schwartz subspace. The same
Hilbertian source-norm argument therefore makes this
restricted map nuclear with target \(\mathscr U_0\);
equivalently one may apply a continuous finite-rank
moment correction to its source nuclear series.
Restricting the return and then applying
the actual theta embedding proves, without a gap
or generic-width condition,
\[
       B_{\phi,g}:K_{a,s,\phi}\longrightarrow J_{a,s}
                      \quad\text{is trace class}
 \tag{MA-MO.830}
\]
and agrees with the original pole-free source action.
There is still no \(K\)-invariance claim or
endomorphism trace for this last map. The previously
proved domain and outgoing-moment obstructions are
unchanged.

\paragraph{What this decides and what it does not decide.}
The theorem applies directly to the requested primitive
nontrivial Dirichlet characters and to both
\(\psi_{269,3}\) sectors with their primitive factors
and structural place set. No use of A, M, or a
generic-width assertion occurs in this return.
It includes widths whose edge contains a zero,
because the input contour has a fixed strict
interior margin and the actual ideal cancels every
zero met during deformation.

The extension is an endomorphism of the completed
theta ideal and has its genuine Fredholm determinant.
It does not construct the full raw common-sum matrix
on \(H\), an inverse on either one-sided Hilbert space,
or a return of discarded quotient factors.
For a general source numerator, the quotient in
(MA-MO.815) has poles; the deformation then has
residues and the original wide inverse need not
be identified by the interior contour. The exterior
and deficient-boundary obstructions MA-MO-024/025
are therefore unaffected.

In particular this result does not show that the
finite-frequency operators of MA-RS-007 converge
on their unrestricted one-sided spaces for every
compact smooth test. Nor does it justify removing
the ideal projection from a common-space trace.
The trace in (MA-MO.803) is compared by the actual
return (MA-MO.801), not assigned from a scalar
zero sum. The full quotient and FT002 arithmetic
deficiency gates remain distinct.

\paragraph{Exact source and framework placement.}
The new path construction uses only the unconditional
analytic inputs and fully proved minimum-modulus
argument in the supplied GRH\_NIGHT MT-M18:
\texttt{work/root/block\_arithmetic\_synthesis.tex},
(M18.5)--(M18.9), together with NT-Z10.
The underlying fixed-Hecke facts have the inherited
primary locators Ramakrishnan and Valenza (1999),
\emph{Fourier Analysis on Number Fields}, \S7.4,
Theorem 7-19, p.~279, and \S7.7, Theorem 7-28,
p.~289. The real-coordinate averaging and the
interior contour estimates are proved here.

The original source inclusion and inverse identification
are MA-MO.173--183, through MT-M18.4 and MT-M24.6--9,
and Meyer (2005), supplied \texttt{main.tex}
lines 3592--3707 and 3808--4043.
The local multiplier, paired source topology and
nuclear comparison have the precise primary
locators recorded in MA-MO-017: Meyer, Lemma 4.3,
Propositions 4.4 and 4.6, Corollary 4.7, split
extension (24), and local trace equation (20);
Tate (1950), Theorem 2.4.1 and \S2.5, thesis
pp.~(2.11), (2.14), (2.16), (2.19).
All finite Fourier phases and conductor/different
monomials are fixed by those actual source maps.

In the inherited cote 1 framework, scans 079--080
and 111--112 through BR-TR-01/02, (MA-MO.821)--(MA-MO.823)
give the required genuine nuclear maps and common
absolutely summable trace pairing. The finite-order
character remains the specified rank-one finite-monodromy
local system, with its Frobenius-determinant Euler
factors. Analytic norm-half centering is not an
integral Tate twist. No new positivity or global
zero bound is inferred from this source return.
